\begin{figure*}[t]
	\centering
	\includegraphics[width=0.9\linewidth]{figs/teaser.pdf}
	\vspace{-4pt}
	\caption{Comparison of various popular architectures to address CL. (a) Traditional dynamic expansion-based CL cannot distinguish unseen data. (b) Zero-shot CL~\cite{zheng2023preventing} suffers from significant computational burdens. (c) The proposed MoE-Adapters and DDAS collaborate to form a parameter-efficient, zero-shot CL.}
	\label{fig:teaser}
 \vspace{-10pt}
\end{figure*}

\section{Introduction}
\label{sec:intro}

Artificial Intelligence (AI), particularly in the realm of large-scale foundation models, has made significant strides in understanding the open world, as evidenced by recent advancements~\cite{OpenAI_2023, touvron2023llama, radford2021learning, liu2023visual, liu2023grounding}. An ideal AI, akin to human cognition, should be able to continuously assimilate new knowledge from the dynamic environment. Traditional fully-supervised training paradigms can't adapt to this scenario due to the high computational costs of integrating new data with historical datasets. In contrast, Continual Learning (CL), offering an efficient incremental training strategy, emerges as a solution by focusing on new data at each training stage. However, CL faces the significant hurdle of ``catastrophic forgetting'' where a model loses previously acquired knowledge upon learning new tasks~\cite{mccloskey1989catastrophic, goodfellow2013empirical}.

To remedy this issue, one of the popular solutions in current CL methods~\cite{aljundi2017expert, donahue2014decaf, girshick2014rich, li2017learning} is to develop dynamic expansion frameworks by incrementally adding task-specific components to a shared base model (see Figure~\ref{fig:teaser} (a)). 
Although these methods show promise in memorization and scalability, they cannot distinguish unseen data and thus overlook zero-shot transfer capability.
Recent advancements like ZSCL~\cite{zheng2023preventing} have brought the zero-shot transfer ability into continual learning by leveraging a pretrained Vision Language Model (VLM). As illustrated in Figure~\ref{fig:teaser} (b), this method relies on knowledge distillation to integrate zero-shot generalization ability from the frozen CLIP and uses parameter regularization to prevent knowledge degradation in continual learning. 
However, these designs often entail large computational burdens and exhibit limitations in long-term memorization. It's then natural to ask whether we can combine the merits of the pretrained foundation model and dynamic expansion strategy to form an effective system with robust memorization and zero-shot transfer abilities.


Recently, Parameter-Efficient Fine-Tune (PEFT) methods \cite{gao2023clip, sung2022vl, zhang2021tip, zaken2021bitfit, hu2021lora} have demonstrated that large-scale models can quickly adapt to downstream tasks via only fine-tuning less-parameterized adapters. This inspires us to build a dynamic expansion framework on VLM with task-specific adapters to relieve the parameter burdens in long-term CL.  Nevertheless, the intuitive approach of stacking adapters during incremental learning introduces a dependency on task identity. This poses challenges in practical scenarios such as class incremental learning where task identity may be unavailable. Furthermore, the use of independent adapters neglects the potential for inter-task knowledge sharing and cooperation, resulting in a limited representation capability and efficacy.


To overcome the outlined challenges, we propose a parameter-efficient continual learning framework by leveraging the recent advance in the field of multi-task learning, Mixture-of-Experts (MoE)~\cite{jacobs1991adaptive, shazeer2017outrageously}. We build a dynamic expansion architecture on a frozen CLIP model~\cite{radford2021learning}, dubbed as incremental MoE-Adapters, in which we take adapters as sparse experts and utilize incrementally incorporated task-specific routers to select the corresponding experts. In the continual learning process, we further apply a novel activate-freeze strategy to help the experts learn intra-task knowledge and encourage inter-task collaboration. Additionally, a Distribution Discriminative Auto-Selector (DDAS) is proposed to automatically allocate the testing data to MoE-Adapters or the pretrained CLIP, enabling effective predictions for seen data and zero-shot transfer for unseen data within a unified framework.




Our extensive experiments across various settings demonstrate the proposed method's effectiveness in addressing the catastrophic forgetting issue, significantly reducing the $60\%$ parameter burdens and memory requirements during training. Furthermore, when applied to few-shot continual learning, the proposed model shows exceptional resistance to forgetting and outperforms the previous arts by $3.6\%$, $7.0\%$ and 4.2\% in a 5-shot setting. 
Our contributions can be summarized as follows:
\begin{itemize}
\item We introduce a parameter-efficient training framework for vision-language models in continual learning, employing a MoE-Adapters based dynamic expansion architecture for enhanced adaptability and efficiency.
\item We develop an incremental activate-freeze strategy in the MoE framework, enabling experts to simultaneously acquire intra-task knowledge and engage in inter-task collaboration.
\item We design a Distribution Discriminative Auto-Selector (DDAS) for automated substream assignment, effectively merging anti-forgetting and zero-shot transfer capabilities within a unified model.
\end{itemize}




